\documentclass[11pt,a4paper]{article}

\usepackage[utf8]{inputenc}
\usepackage[T1]{fontenc}
\usepackage[margin=2.2cm]{geometry}
\usepackage{parskip}
\usepackage{titlesec}
\usepackage{enumitem}
\usepackage{booktabs}
\usepackage{xcolor}
\usepackage{tcolorbox}
\usepackage{hyperref}
\usepackage{fancyhdr}
\usepackage{listings}
\usepackage{longtable}
\usepackage{array}

\tcbuselibrary{skins,breakable}

\definecolor{accent}{RGB}{30,80,150}
\definecolor{softgray}{RGB}{245,246,248}
\definecolor{codebg}{RGB}{248,248,250}
\definecolor{warn}{RGB}{160,60,40}
\definecolor{good}{RGB}{30,120,70}

\titleformat{\section}{\Large\bfseries\color{accent}}{\thesection}{0.8em}{}
\titleformat{\subsection}{\large\bfseries}{\thesubsection}{0.6em}{}
\titleformat{\subsubsection}{\normalsize\bfseries\itshape}{\thesubsubsection}{0.5em}{}

\hypersetup{
  colorlinks=true, linkcolor=accent, urlcolor=accent,
  pdftitle={AgentIA --- Measured Reliability, Cost and Readiness},
  pdfauthor={AgentIA Engineering}
}

\pagestyle{fancy}
\fancyhf{}
\fancyhead[L]{\small\color{accent}AgentIA --- Measurement and Readiness}
\fancyhead[R]{\small\color{accent}\thepage}
\renewcommand{\headrulewidth}{0.4pt}

\lstset{
  basicstyle=\ttfamily\small, backgroundcolor=\color{codebg},
  frame=single, rulecolor=\color{softgray}, breaklines=true,
  columns=fullflexible, keepspaces=true, showstringspaces=false,
  xleftmargin=0.5em, xrightmargin=0.5em,
}

\newtcolorbox{notebox}[1][]{colback=softgray,colframe=accent,boxrule=0.4pt,
  arc=2pt,left=6pt,right=6pt,top=4pt,bottom=4pt,breakable,#1}

\title{\vspace{-1.4cm}\textbf{Measured Reliability, Cost and Readiness}\\[0.25cm]
\large What the platform actually does on two corpora, what it costs per task, and
what stands between it and production}
\author{AgentIA Engineering}
\date{2026-09-29}

\begin{document}
\maketitle
\thispagestyle{fancy}

\begin{notebox}
\textbf{What this document is.} The measured state, replacing the estimates in
\texttt{agentia\_first\_baseline.tex} (whose central claim has been retracted
in-document) and the stale figures in \texttt{agentia\_current\_state.tex}.

Every number here comes from a real run against a real provider with real container
verification attempted. Costs are read from the platform's own store
(\texttt{backend/cost\_tracking.db}) and independently confirmed in MLflow. Where a
figure could not be measured, this document says so rather than estimating it.

\textbf{The headline is a null result, and it is the most useful thing here.} A task
corpus made ten times harder changed nothing: the generator satisfies every
encodable rule on its first attempt, every time. That closes the optimization line
structurally, and it redirects the effort to the one thing still unmeasured.
\end{notebox}

\tableofcontents
\vspace{0.5cm}

\section{What was measured}

Two corpora, run end to end on 2026-09-29, provider \texttt{deepseek}, model
\texttt{deepseek-flash}:

\begin{center}
\begin{tabular}{@{}p{0.22\textwidth}p{0.16\textwidth}p{0.54\textwidth}@{}}
\toprule
\textbf{Corpus} & \textbf{Size} & \textbf{Character} \\
\midrule
frozen five & 5 blueprints & the held-out set frozen by feature 011; deliberately
easy \\
hard ten & 10 blueprints & 53 entities, 356 attributes, 80 stories, 207 scenarios;
each stresses a different difficulty axis (referential depth, constraint variety,
naming collisions, layering temptation, paging, mutation lifecycle, batch
scheduling, breadth, invariant-heavy unhappy paths, heterogeneous identifiers) \\
\bottomrule
\end{tabular}
\end{center}

\section{Result: the corpus was not the problem}

\begin{center}
\begin{tabular}{@{}lrr@{}}
\toprule
\textbf{Measure} & \textbf{frozen five} & \textbf{hard ten} \\
\midrule
Sessions recorded & 5 & 10 \\
Sessions that generated & \textbf{5 / 5} & \textbf{10 / 10} \\
Stage runs succeeding on the FIRST request & 25 / 25 & \textbf{50 / 50} \\
Correction attempts consumed & 0 & \textbf{0} \\
Artifacts per session & 15--23 & \textbf{39--63} \\
Conformance score & 100 $\times$ 5 & \textbf{100 $\times$ 10} \\
Conformance density & 0.0 $\times$ 5 & \textbf{0.0 $\times$ 10} \\
Findings in the final artifact sets & none & \textbf{none} \\
\bottomrule
\end{tabular}
\end{center}

A corpus with ten times the structural complexity produced \emph{identical}
conformance. The hypothesis ``the tasks are too easy'' is falsified.

\subsection{Why the measure cannot vary, at any corpus size}
\label{sec:structural}

This is structural, not statistical, and it is the finding that closes the
optimization line.

The compliance gate **blocks stage-local violations before persistence**. An artifact
set that reaches the workspace has therefore already passed every local rule by
construction. The only findings that can survive into a saved set are
\emph{accumulated} whole-project ones --- and exactly one exists
(\texttt{PRINCIPLE\_III\_CENTRALIZED\_ERRORS}, the requirement for a
\texttt{@RestControllerAdvice}), which the generator satisfies by emitting
\texttt{GlobalExceptionHandler} in the controller stage.

So the final-set conformance measure is pinned at 100 for **any session that
generates at all**. The feature's own spec anticipated this and could not quantify
it:

\begin{notebox}
``because stage-local violations are \textbf{blocked before persistence}, the
findings that survive into a saved artifact set come from whole-project rules only,
which is a \textbf{narrow} signal whose real-world variance is unmeasured.''
\end{notebox}

It is now measured. Variance is \textbf{zero}, at both levels of the signal:

\begin{center}
\begin{tabular}{@{}p{0.30\textwidth}p{0.62\textwidth}@{}}
\toprule
\textbf{Signal level} & \textbf{Observed across all 15 sessions} \\
\midrule
Final artifact set & 100 / density 0.0 / zero findings, in every session \\
Transient (per-stage introduction) & \emph{identical} in every session:
\texttt{PRINCIPLE\_III\_CENTRALIZED\_ERRORS} $\times$ 3, the accumulated whole-project
rule firing during the three stages that precede the controller \\
\bottomrule
\end{tabular}
\end{center}

\textbf{Consequence for optimization.} An objective with no variance has no
headroom. No corpus growth and no task difficulty can create it, because the gate
removes the local signal before it is stored. Reflective skill optimization is
therefore closed \emph{structurally}, not for want of samples.

\section{Result: cost and tokens per task}

Generation now costs one request per stage --- five per session --- because nothing
fails validation and needs a second attempt. Cost scales with \emph{blueprint
complexity} (payload and output tokens), never with retries.

\subsection{Hard ten --- per task}

\begin{center}
\begin{tabular}{@{}lrrrr@{}}
\toprule
\textbf{Task} & \textbf{Calls} & \textbf{Input tok} & \textbf{Output tok} & \textbf{Cost USD} \\
\midrule
hard-01 & 5 & 74{,}649 & 107{,}044 & \$0.0745 \\
hard-02 & 5 & 55{,}457 & 77{,}352 & \$0.0538 \\
hard-03 & 5 & 61{,}780 & 91{,}533 & \$0.0633 \\
hard-04 & 5 & 49{,}470 & 75{,}466 & \$0.0518 \\
hard-05 & 5 & 49{,}827 & 88{,}788 & \$0.0598 \\
hard-06 & 5 & 57{,}412 & 83{,}082 & \$0.0575 \\
hard-07 & 5 & 65{,}163 & 98{,}729 & \$0.0681 \\
hard-08 & 5 & 78{,}794 & 110{,}118 & \$0.0770 \\
hard-09 & 5 & 73{,}428 & 84{,}419 & \$0.0607 \\
hard-10 & 5 & 61{,}887 & 101{,}287 & \$0.0691 \\
\midrule
\textbf{Total} & \textbf{50} & \textbf{627{,}867} & \textbf{917{,}818} & \textbf{\$0.6357} \\
\bottomrule
\end{tabular}
\end{center}

\subsection{Per-stage breakdown (one simple service)}

Every stage is one request; the test stage is the most expensive because its output
is the largest.

\begin{center}
\begin{tabular}{@{}lrrr@{}}
\toprule
\textbf{Stage} & \textbf{Input tok} & \textbf{Output tok} & \textbf{Cost USD} \\
\midrule
SCAFFOLDER & 2{,}200 & 6{,}620 & \$0.0040 \\
DOMAIN & 2{,}669 & 1{,}759 & \$0.0012 \\
SERVICE & 3{,}250 & 5{,}166 & \$0.0034 \\
CONTROLLER & 3{,}925 & 4{,}210 & \$0.0029 \\
TEST & 5{,}711 & 9{,}108 & \$0.0061 \\
\midrule
\textbf{Total} & \textbf{17{,}755} & \textbf{26{,}863} & \textbf{\$0.0176} \\
\bottomrule
\end{tabular}
\end{center}

\subsection{Pricing summary}

\begin{center}
\begin{tabular}{@{}lrrr@{}}
\toprule
\textbf{Service size} & \textbf{Artifacts} & \textbf{Tokens (in / out)} & \textbf{Cost} \\
\midrule
simple (\texttt{minimal}) & 15 & 17{,}755 / 26{,}863 & \textbf{\$0.0176} \\
medium (frozen five range) & 15--23 & $\approx$25k / $\approx$60k & \$0.018--0.025 \\
hard (hard ten range) & 39--63 & 49k--79k / 75k--110k & \textbf{\$0.052--0.077} \\
\midrule
per model call (all runs) & --- & --- & $\approx$ \textbf{\$0.005} \\
\bottomrule
\end{tabular}
\end{center}

\subsection{Telemetry: the mirror works, and it agrees}

Cost records are written to the local SQLite store (the system of record) and
mirrored per call \emph{and} per session to MLflow. Reading the hard ten back from
the tracking server independently reproduced the local totals exactly --- 50 runs,
\$0.6357 --- which is what makes the local figure trustworthy rather than merely
self-consistent.

\begin{lstlisting}
# authoritative, offline, deterministic
.venv/bin/python backend/scripts/report_task_costs.py --prefix hard-r1
# the mirror, for a UI
.venv/bin/mlflow server --backend-store-uri sqlite:///mlflow.db \
    --default-artifact-root ./mlruns --host 127.0.0.1 --port 5000
\end{lstlisting}

\begin{notebox}[colframe=warn]
\textbf{Two defects had to be fixed before these numbers meant anything.} Records and
cost rows were keyed differently (a fresh session prefix disambiguated the record but
not the cost, so a run that spent \$0.1235 reported ``cost: no data''), and session
aggregates were never mirrored at all --- only per-call records. Both are fixed, and
the record/cost identity is now pinned by a contract test. Cost telemetry that
disagrees with itself is worse than none.
\end{notebox}

\section{How this differs from a developer using generative AI}

The honest answer first: \textbf{the model is the same commodity model a developer
would call.} There is no fine-tuning, no custom weights, and no claim to better code.
The difference is not the generator; it is the governed pipeline around it.

\begin{center}
\begin{tabular}{@{}p{0.20\textwidth}p{0.34\textwidth}p{0.38\textwidth}@{}}
\toprule
& \textbf{Developer + generic AI} & \textbf{AgentIA} \\
\midrule
Input & a prompt, refined by conversation & a formal blueprint: entities,
attributes, validation rules, BDD acceptance scenarios, frozen before generation \\
Verification & the developer's own, by hand, at their discretion & a
\textbf{deterministic rule gate} plus a \textbf{hermetic offline build}
(\texttt{mvn test -o}, no network) --- executable, repeatable, not a judgement call \\
Failure record & chat history, discarded & a durable per-session record: stage,
outcome, requests spent, rule, artifact, severity, and which stage introduced it \\
Honesty about verification & unknown whether the code was ever run & \textbf{refuses
to substitute a synthetic success}; a session that could not be verified is marked
and excluded from every figure \\
Reproducibility & not reproducible & a deterministic mode that is byte-identical for
identical input; model runs carry provider, model, instruction revision and content
digests \\
Cost & a seat licence or an untracked token bill & priced per call, per stage, per
session, mirrored, and reportable per task \\
Self-correction & manual re-prompting, unbounded & bounded at three attempts per
stage, recorded, and it stops rather than looping \\
\bottomrule
\end{tabular}
\end{center}

\subsection{And the honest counterpoints}

\begin{itemize}[leftmargin=1.4em]
  \item \textbf{The gate never had to fire in these runs.} The model passed every
  rule first time. The gate's value is insurance, and insurance looks worthless
  right up until it does not. Nothing measured here demonstrates that the gate
  catches errors, because there were none to catch.
  \item \textbf{Conformance is not correctness.} The instrument is a fixed,
  hand-written rule set. A well-layered but broken service scores full marks.
  \item \textbf{The build and test step is unmeasured.} 0 of 15 sessions were
  verified: the container runtime was unreachable from the measurement shell. Every
  session terminated at the sandbox with an honest ``could not verify''. Whether the
  generated code compiles and passes its tests is therefore \emph{unknown}, and it is
  the single most important unknown in this document.
  \item \textbf{One product entry point does not use the model.}
  \texttt{POST /sessions} runs deterministic templates while quick-start and the
  orchestrator routes run the model (\S\ref{sec:blocker}). The measurements above
  describe the model path.
\end{itemize}

\section{Readiness: what stands between this and production}
\label{sec:blocker}

\begin{notebox}[colframe=warn]
\textbf{One of the two generation entry points silently does not use the model.}
The platform has two ways to start a session, and they disagree:

\begin{itemize}[leftmargin=1.4em]
  \item \textbf{Sequential path.} \texttt{POST /sessions/quick-start} with
  \texttt{autoRun: true}, and \texttt{/orchestrator/pipeline/run}, thread the caller's
  credentials into \texttt{pipeline\_runner}, which calls
  \texttt{select\_generation\_mode(api\_key, provider, model\_name)} and selects
  \textbf{MODEL}. \textbf{These routes can use the model}, and they are the ones a
  client should use today.
  \item \textbf{Graph path.} \texttt{POST /sessions} builds its state by calling
  \texttt{select\_generation\_mode()} \emph{with no arguments} and passes no
  \texttt{llm\_api\_key}. With no key and no provider that returns
  \textbf{DETERMINISTIC}, so this route emits offline templates and never calls the
  model --- although the graph itself supports MODEL mode, as the batch driver proves
  by invoking the same graph with credentials in the state.
\end{itemize}

The consequence is a product surface that behaves differently depending on which
endpoint a client picks: one returns model-generated code, the other templates, and
nothing in the response distinguishes them. It is a one-line wiring fix, not a
blocker --- and it is a correction to an earlier draft of this document, which claimed
the application could not reach the model at all. That was wrong: the quick-start and
orchestrator routes do.
\end{notebox}

Ranked by what unblocks the most:

\begin{center}
\begin{tabular}{@{}p{0.07\textwidth}p{0.30\textwidth}p{0.55\textwidth}@{}}
\toprule
\textbf{Priority} & \textbf{Item} & \textbf{Why it is where it is} \\
\midrule
P0 & Entry-point consistency & \texttt{POST /sessions} silently runs offline
templates while quick-start runs the model; a client cannot tell which it received \\
P1 & \textbf{Build pass rate} & the only remaining measurement. Run the corpus where
the container works. Until then ``reliable'' is unfalsifiable \\
P2 & Fix build failures in \emph{code} & if builds fail, the diagnostic channel names
the classes; systematic fixes belong in the harness, the instruction set or the
dependency allowlist --- deterministic, reviewable, testable --- not in prose skills
edited by a model \\
P3 & Verification-guided selection (best-of-$k$) & the evidence-ranked lever. The
verifier already exists and is free; generation is \$0.064/task, so $k{=}4$ is
$\approx$\$0.26 and converts unreliable generation into reliable output \emph{without
trusting the model} \\
P4 & Operational hardening & container runtime reachable where it is deployed;
per-session cost attribution; the repair cap (the constitution says 3,
\texttt{config.py} says 5); the stale README; the two empty skill placeholders ---
author or delete them \\
\bottomrule
\end{tabular}
\end{center}

\section{Recommendation}

\begin{enumerate}[leftmargin=1.4em]
  \item \textbf{Stop the skill-optimization line.} Not for cost --- it is
  \$0.06/session --- but because the objective is structurally degenerate
  (\S\ref{sec:structural}). Record it as closed on evidence, keep the components on
  the shelf, and do not grow the corpus for it.
  \item \textbf{Make the two entry points agree} (P0). Quick-start reaches the
  model; \texttt{POST /sessions} silently does not. A client should not have to know
  which endpoint yields model output. It is a small wiring fix.
  \item \textbf{Measure the build pass rate} (P1). One run in a shell where the
  container runtime initialises. No further reliability work is defensible before it:
  everything else is optimising a pipeline whose output has never been proven to
  compile.
  \item \textbf{If builds fail, fix them in code} (P2). The lesson of this whole
  exercise is that the reliable place to encode a rule is a deterministic check and a
  test --- three real defects were found this way and none by the language model.
  \item \textbf{Then, and only then, spend on best-of-$k$} (P3). It is measurable at
  this corpus size because it moves a binary per-task outcome, where a conformance
  nudge of a few points would be buried in noise.
  \item \textbf{Do not} present conformance as a reliability metric, and do not
  resurrect the optimizer without first exhibiting a measure that varies.
\end{enumerate}

\subsection*{The one-sentence position}

The generator is reliable on the axis that was measured --- first-attempt
rule-compliance, at \$0.064 per hard service --- but the axis that matters for
production, whether the output builds and passes its tests, has not been measured
once; and one of the application's two generation entry points never invokes the model
at all.

\end{document}
